\documentclass[11pt,a4paper]{article}

\usepackage[utf8]{inputenc}
\usepackage[T1]{fontenc}
\usepackage{lmodern}
\usepackage[margin=0.78in]{geometry}
\usepackage[table]{xcolor}
\usepackage{array}
\usepackage{tabularx}
\usepackage{booktabs}
\usepackage{enumitem}
\usepackage{titlesec}
\usepackage{fancyhdr}
\usepackage{lastpage}
\usepackage{hyperref}

\definecolor{TextInk}{HTML}{233142}
\definecolor{HeadingBlue}{HTML}{1D3557}
\definecolor{AccentTeal}{HTML}{2A7F92}
\definecolor{AccentGold}{HTML}{C58B4E}
\definecolor{TableStripe}{HTML}{F3F8FA}
\definecolor{SoftBlue}{HTML}{EEF6F8}

\hypersetup{
  colorlinks=true,
  linkcolor=HeadingBlue,
  urlcolor=AccentTeal,
  pdftitle={Geography Paper 1 SL Set 2 Student Work Marking Report},
  pdfauthor={IB Geography SL Preparation}
}

\color{TextInk}
\setlength{\parskip}{0.45em}
\setlength{\parindent}{0pt}
\setlength{\tabcolsep}{5pt}
\renewcommand{\arraystretch}{1.18}
\setlist[itemize]{leftmargin=1.35em,itemsep=3pt,topsep=4pt}

\newcolumntype{Y}{>{\raggedright\arraybackslash}X}
\newcolumntype{L}[1]{>{\raggedright\arraybackslash}p{#1}}
\newcommand{\term}[1]{\textcolor{HeadingBlue}{\textbf{#1}}}
\newcommand{\score}[1]{\textcolor{AccentTeal}{\textbf{#1}}}
\newcommand{\boxnote}[3]{%
\noindent\colorbox{#1}{%
  \begin{minipage}{0.965\textwidth}
  \sffamily
  \textbf{\textcolor{HeadingBlue}{#2}} #3
  \end{minipage}
}}

\titleformat{\section}
  {\normalfont\Large\sffamily\bfseries\color{HeadingBlue}}
  {}{0pt}{}
  [\vspace{0.08em}{\color{AccentGold}\titlerule[0.8pt]}]

\titleformat{\subsection}
  {\normalfont\large\sffamily\bfseries\color{AccentTeal}}
  {}{0pt}{}

\pagestyle{fancy}
\setlength{\headheight}{14pt}
\fancyhf{}
\fancyhead[L]{\sffamily\color{AccentTeal}Teacher Marking Report}
\fancyhead[R]{\sffamily\color{HeadingBlue}May 2026 Paper 1 Set 2}
\fancyfoot[C]{\sffamily\color{HeadingBlue}\thepage\ of \pageref{LastPage}}
\renewcommand{\headrulewidth}{0.4pt}

\begin{document}

\begin{center}
  {\sffamily\bfseries\color{HeadingBlue}\LARGE Geography Paper 1 SL Set 2 Marking Report}\\[0.25em]
  {\sffamily\color{AccentTeal}\large Options F and G Student Work}
\end{center}

\boxnote{SoftBlue}{Basis for marking.}{This report marks the Set 2 student
work markdown against the Set 2 question paper and marking guide in
\path{plans/}. The student answered Option F, Questions 11 and 12(a), and
Option G, Questions 13 and 14(a). Spelling and grammar are ignored unless they
affect geographical meaning.}

\section{Overall Mark}

\rowcolors{2}{TableStripe}{white}
\begin{tabularx}{\textwidth}{L{0.16\textwidth}L{0.14\textwidth}Y}
\toprule
\textbf{Question} & \textbf{Mark} & \textbf{Teacher rationale} \\
\midrule
Q11 & \score{9/10} & The short explanations are strong and mostly follow clear
cause-effect chains. The only lost mark is for 11(a)(i), where the student
identifies the irrigated core instead of the dryland smallholder region as
having the highest percentage of food-insecure households. \\
Q12(a) & \score{6/10} & Reaches the 5--6 band. The answer uses the Democratic
Republic of Congo and explains shared drivers such as poor infrastructure,
food insecurity, sanitation and healthcare access, plus the effect of
undernutrition on disease risk. It remains below the 7--8 band because the
two-way relationship is incomplete and the place evidence is general. \\
Q13 & \score{9/10} & Most short answers meet the markscheme. The student loses
one mark for 13(a)(i), where the central core is given instead of the informal
hillside zone for the highest share of walking or cycling trips. \\
Q14(a) & \score{8/10} & A secure 7--8 band essay. New York City is used as a
named example, and congestion charging/tolls and public transport investment
are examined with strengths and limitations. It does not reach 9--10 because
the comparison of strategies and evidence of effectiveness are not systematic
enough for the top band. \\
\midrule
\textbf{Total} & \score{32/40} & Strong performance on developed short-answer
questions and a good traffic-congestion essay. The main lost marks come from
two incorrect graph readings and from an extended response on malnutrition and
infectious disease that does not fully examine the reciprocal links. \\
\bottomrule
\end{tabularx}
\rowcolors{2}{}{}

\section{Option F: Food and Health}

\subsection{Q11: Dry Season, Food Prices and Water Access \score{9/10}}

\rowcolors{2}{TableStripe}{white}
\begin{tabularx}{\textwidth}{L{0.16\textwidth}L{0.12\textwidth}Y}
\toprule
\textbf{Part} & \textbf{Mark} & \textbf{Rationale} \\
\midrule
11(a)(i) & \score{0/1} & The correct answer is \term{Dryland smallholder}, or
46\%. The student gives \term{Irrigated core}, which has a much lower
food-insecure-household value on the graph. \\
11(a)(ii) & \score{1/1} & \term{21\%} is the correct percentage of households
reporting child wasting in the dryland smallholder region. \\
11(b) & \score{2/2} & The response identifies a valid mechanism: a dry season
reduces rainfall and forces rural users to rely more heavily on stored water
or groundwater. It develops this by linking agricultural water demand to lower
remaining access to safe water. \\
11(c)(i) & \score{3/3} & The answer gives a clear chain from a climate shock to
reduced crop productivity and yield, then to undernutrition and child wasting.
The explanation is directly relevant to vulnerable rural households. \\
11(c)(ii) & \score{3/3} & The answer explains that lower crop yields reduce
food supply, while food remains a necessity, increasing demand pressure and
raising food prices in urban low-income districts. \\
\bottomrule
\end{tabularx}
\rowcolors{2}{}{}

\subsection{Q12(a): Malnutrition and Infectious Disease \score{6/10}}

This response reaches the upper part of the 5--6 markband. It is relevant,
structured and uses a named place, the Democratic Republic of Congo. The
student understands that malnutrition and infectious disease often occur in
the same places because of shared contextual factors, especially inadequate
food storage and distribution, weak infrastructure, poor sanitation, and
uneven access to healthcare between Kinshasa and rural areas.

The answer also explains one important biological link: undernutrition reduces
the body's ability to stay healthy, making people more vulnerable to
infectious disease. This is creditworthy because it links food access,
nutrient intake, immune function and disease risk.

The response is held at \score{6/10} because the link is not examined fully in
both directions. A stronger answer would also explain how infectious disease
can worsen malnutrition by reducing appetite, increasing nutrient loss, causing
diarrhoea, reducing absorption, and slowing child growth. The DRC example is
useful, but it is broad; there is limited disease-specific evidence such as
cholera, diarrhoeal disease, measles, malaria, TB or other named infectious
diseases. The conclusion is relevant but mostly repeats the earlier argument
rather than weighing the relative importance of shared causes and direct
biological feedbacks.

To move this response into the \score{7--8} band, the student should develop a
clear two-way cycle: malnutrition weakens immunity and increases infection
risk, while infectious disease worsens nutrient loss and food insecurity. To
reach \score{9--10}, this should be supported with more precise place evidence
and a balanced judgment about whether shared poverty, WASH conditions and
health systems are more important than the direct biological feedbacks.

\section{Option G: Urban Environments}

\subsection{Q13: Urban Transport Patterns \score{9/10}}

\rowcolors{2}{TableStripe}{white}
\begin{tabularx}{\textwidth}{L{0.16\textwidth}L{0.12\textwidth}Y}
\toprule
\textbf{Part} & \textbf{Mark} & \textbf{Rationale} \\
\midrule
13(a)(i) & \score{0/1} & The correct answer is \term{Informal hillside}, or
45\%, for the highest share of walking or cycling trips. The student gives
\term{Central core}. \\
13(a)(ii) & \score{1/1} & \term{Outer suburb} is correct for the zone with the
highest share of private-car trips. \\
13(b) & \score{2/2} & The response identifies high greenhouse-gas emissions
from car dependence and develops this through high household car ownership in
outer urban zones. This is a valid environmental-stress chain. \\
13(c)(i) & \score{3/3} & The student uses employment and business
centralization as an economic factor. The answer links job location to travel
distance and then to transport choice, which is enough development for full
credit. \\
13(c)(ii) & \score{3/3} & The response identifies transport technology and
infrastructure, such as light rail or subway systems, and explains that faster
commute times can shift people towards public transport. \\
\bottomrule
\end{tabularx}
\rowcolors{2}{}{}

\subsection{Q14(a): Managing Traffic Congestion in New York City \score{8/10}}

This essay securely reaches the 7--8 markband. It is focused on the command
term \term{examine} and uses New York City as a named city throughout. The
student explains why congestion matters, linking idling traffic to emissions
and lost economic productivity, then examines two relevant strategies:
tolls/congestion charging and investment in public transportation.

The tolls paragraph is the stronger part of the essay. It explains the demand
management logic: charging vehicles entering congested areas discourages car
use, reduces the number of vehicles in the toll zone and can reduce travel
times. The student also evaluates the strategy by noting financial burdens on
drivers and possible displacement of congestion to alternative routes.

The public transport paragraph is also relevant. It explains that toll revenue
can be invested in the MTA and that an extensive subway system can reduce the
number of cars on the road. The use of the 4.6 million daily subway users is a
useful piece of case-specific support. The weakness about uneven subway access
is valid because public transport only reduces congestion effectively where it
is accessible and reliable.

The essay does not reach \score{9--10} because the evaluation is not fully
systematic. A top-band answer would compare the two strategies more directly,
use more precise evidence for reductions in vehicle numbers or travel speeds,
and consider time scale, equity and spatial impacts in more depth. It would
also distinguish more carefully between managing congestion itself and reducing
air pollution or greenhouse-gas emissions as related outcomes.

\section{Main Targets}

\begin{itemize}
  \item Read graph categories carefully before answering 1-mark data-response
  questions; two marks were lost from incorrect graph identification.
  \item In 10-mark essays, examine both sides of a relationship where the
  question asks for \term{links}, especially feedback cycles such as
  malnutrition causing disease vulnerability and disease worsening
  malnutrition.
  \item Add more precise case-study evidence in essays: named diseases,
  statistics, specific locations within the country or city, and measured
  outcomes of strategies.
  \item For \term{examine} questions, keep the useful strength/weakness
  structure but make the final judgment more comparative and evidence-based.
\end{itemize}

\end{document}
